\documentclass[xcolor=pdftex,dvipsnames,table,mathserif]{beamer}
\usepackage{subfigure}
\usepackage{amsbsy}
\usepackage{tikz}
\usetikzlibrary{arrows}
\usepackage{amsmath,graphicx,dsfont,color}
\usepackage{amsfonts}
\usepackage{amssymb}
\usepackage{array}

\bibliographystyle{apalike}

\setbeamertemplate{bibliography item}{\insertbiblabel}
\setbeamertemplate{bibliography entry title}{}
\setbeamertemplate{bibliography entry location}{}
\setbeamertemplate{bibliography entry note}{}

\DeclareMathOperator*{\argmin}{arg\,min}
\DeclareMathOperator*{\argmax}{arg\,max}
%Definitiona

\newcommand{\x}{\mathbf{x}}
\newcommand{\X}{\mathbf{X}}
\newcommand{\W}{\mathbf{W}} %Weight
\newcommand{\bais}{\mathbf{b}}%Bais
\newcommand{\act}{\texttt{g}}%Activation
\newcommand{\loss}{L}
\newcommand{\pdata}{\hat{p}_{\texttt{data}}}
\newcommand{\nsize}{n}
\newcommand{\param}{\boldsymbol{\theta}}
\newcommand{\featmap}{\boldsymbol{\phi}}
\newcommand{\EV}{\mathbb{E}}







\usepackage{physics}

\graphicspath{{../graphics/}}

%% \usepackage{animate}

\AtBeginSection[]{
  \begin{frame}{Contents}
  \tableofcontents[currentsection, hideothersubsections]
  \end{frame}
}

\AtBeginSubsection[]{
  \begin{frame}{Contents}
  \tableofcontents[currentsection, subsectionstyle=show/shaded/hide]
  \end{frame}
}

\setbeamertemplate{footline}[frame number]{}
\setbeamertemplate{navigation symbols}{}
\setbeamertemplate{section in toc}[square]
\setbeamertemplate{items}[square]

\title{Deep Learning for Image Analysis\\Course Introduction}
\author{E. Decencière, Thomas Walter, Santiago Velasco-Forero}
\date{MINES ParisTech\\
  PSL Research University
}
\titlegraphic{\includegraphics[height=1.7cm]{../graphics/logoemp}}

\useinnertheme{rounded}
\usecolortheme{rose}

%%%%%%%%%%%%%%%%%%%%%%%%%%%%%%%%%%%%%%%%%%%%%%%%%%
%%%%%%%%%%%%%%%%%%%%%%%%%%%%%%%%%%%%%%%%%%%%%%%%%%
\begin{document}
\begin{frame}
\titlepage
\end{frame}


%%%%%%%%%%%%%%%%%%%%%%%%%%%%%%%%%%%%%%%%%%%%%%%%%%
\frame{
  \frametitle{About the lecturers}

\begin{columns}
  \begin{column}{.2\textwidth}
\vfill
\includegraphics[width=\textwidth]{Thomas_Walter_Heidelberg.jpg}\\
\vspace{2em}
    \includegraphics[width=\textwidth]{velascoforero}\\
\vspace{2em}
    \includegraphics[width=\textwidth]{ed.jpg}

  \end{column}
  \begin{column}{.8\textwidth}

    \begin{block}{Thomas Walter \hfill \scriptsize{\url{http://members.cbio.mines-paristech.fr/\~twalter}}}
      \scriptsize{
    \begin{itemize}
    \item Research: bioimage informatics
    \item Main application fields: High Content Screening (HCS), as a method to systematically study biological processes by analyzing cellular phenotypes
    \end{itemize}
    }
  \end{block}

  \begin{block}{Santiago Velasco-Forero  \hfill \scriptsize{\url{http://cmm.mines-paristech.fr/\~velasco}}}
      \scriptsize{
    \begin{itemize}
    \item Research: image processing, pattern recognition, multivariate statistics, graph-based data/image analysis
    \item Main application fields: Remote Sensing, cosmetology, astronomy, hyperspectral imaging.
    \end{itemize}
    }
  \end{block}

    \begin{block}{Etienne Decencière \hfill \scriptsize{\url{http://cmm.mines-paristech.fr/\~decenciere}}}
      \scriptsize{
    \begin{itemize}
    \item Research: mathematical morphology and image analysis
    \item Main application fields: Ophthalmology, dermatology, cosmetology, astronomy
    \end{itemize}
    }
  \end{block}

  \end{column}
\end{columns}

}

%%%%%%%%%%%%%%%%%%%%%%%%%%%%%%%%%%%%
\begin{frame}{Objectif et pré-requis}

\begin{block}{Objectif}
  Introduction à la théorie et à la pratique de l'apprentissage profond pour l'analyse d'images.
\end{block}

\begin{block}{Connaissances préalables}
\begin{itemize}
\item Algèbre et probabilités de base
\item Programmation: python
\end{itemize}
  \end{block}

  \begin{block}{Langue}
    \begin{itemize}
    \item Slides: anglais
    \item Oral: anglais ou français (en fonction de l'auditoire)
    \end{itemize}

\end{block}


\end{frame}



%%%%%%%%%%%%%%%%%%%%%%%%%%%%%%%%%%%%
\begin{frame}{Contenu}


\begin{block}{Approche pédagogique}

  \begin{itemize}
  \item 8 demi-journées:
    \begin{itemize}
    \item Cours magistraux
    \item Conférences thématiques par des invités extérieurs
    \end{itemize}
  \item Jupyter notebooks (python, keras) à travailler en autonomie
  \end{itemize}
\end{block}

\begin{block}{Sujets abordés}
\begin{itemize}
\item De la classification à la transformation d'images
\item Introspection des réseaux de neurones
\item Supervision faible
\item Autoencodeurs et réseaux antagonistes génératifs (GANs)
\item Apprentissage de métrique
\end{itemize}
\end{block}

  \begin{block}{Evaluation}
  Examen
  \end{block}

\end{frame}


%%%%%%%%%%%%%%%%%%%%%%%%%%%%%%%%%%%%%%%%%%%%%%%%%%
%%%%%%%%%%%%%%%%%%%%%%%%%%%%%%%%%%%%%%%%%%%%%%%%%%

\end{document}
